% Principal-cut continuation of the preceding translation identities.
\section{Transport of Stieltjes translation integrals across the cut}
\label{stcut:section}
The positive-shift identities admit a contour prescription beyond the
negative real axis. For the digamma reconstruction the prescription
requires only Cauchy principal values and finitely many residues.
For integrals of higher Stieltjes functions themselves, logarithmic
cut terms must be retained as boundary distributions. This gives a
concrete answer to the contour-transport question in
\cite[Further research, question 2]{stmell:balanced}, within the
translation Mellin and reconstruction families considered here.
The classical inputs are the Hurwitz shift relations
\cite[Eqs.~25.11.3--25.11.4]{intdist:DLMFHurwitz}, their zeroth
Laurent coefficient (the digamma recurrence), and elementary
Cauchy boundary values. No historical-priority claim is intended.

\subsection{Complex shifts before reaching the cut}
Let $\Omega=\mathbb C\setminus(-\infty,0]$, with the principal logarithm.
Continue $\zeta(s,a)$ and $\gamma_n(a)$ in $a\in\Omega$ using the
Hurwitz shift relation. Let $a_1,\ldots,a_J\in\Omega$ and
$\sum_jc_j=0$. Define
\[
 f(x)=-\sum_jc_j\psi(x+a_j),\qquad d_n=\sum_jc_j\gamma_n(a_j).
\]
Then the digamma-only reconstruction \eqref{stmell:digamma-reconstruct}
holds unchanged as an absolutely convergent complex integral. So does
its Mellin generator
\begin{equation}\label{stcut:complex-master}
 \int_0^\infty x^{v-1}f(x)\,dx
  =\frac{\pi}{\sin\pi v}\sum_jc_j\zeta(1-v,a_j),
 \qquad0<\operatorname{Re}v<1.
\end{equation}
Here $x^{v-1}$ uses the real logarithm of positive $x$.
Indeed the integrands depend holomorphically on every $a_j$ on
$\Omega^J$. On compact subsets they are bounded at zero and are
$O(x^{-1})$ at infinity after the zero-sum cancellation, with the
same bounds for fixed parameter derivatives. The corresponding
subtracted logarithmic integrands are also uniformly integrable.
The identity theorem, applied successively in the shift parameters,
extends the positive-shift formulas to $\Omega^J$.

\subsection{Principal-value reconstruction on the two banks}
Fix real $a_j\notin\{0,-1,-2,\ldots\}$ and complex weights of sum
zero. For $\sigma\in\{+1,-1\}$, define
\[
 d_n^\sigma=\lim_{\eta\downarrow0}\sum_jc_j\gamma_n(a_j+\sigma i\eta).
\]
The limit uses the upper bank for $\sigma=+1$. Put
\[
 \mathcal B=\{b>0:b=-a_j-k\text{ for some }j,\ k\in\mathbb Z_{\ge0}\},
 \qquad R_b=\sum_{\substack{1\le j\le J,\ k\in\mathbb Z_{\ge0}\\a_j+k=-b}}c_j.
\]
In the sum defining $R_b$, $j$ ranges over $1,\ldots,J$ and $k$ over
nonnegative integers. Equal pole locations are combined. Thus
$\mathcal B$ is finite and $R_b$ is exactly the residue of
$f(x)=-\sum_jc_j\psi(x+a_j)$ at $x=b$. A zero residue removes that
pole. The common value $d_0^+=d_0^-=f(0)$ is denoted by $d_0$.

For $n\ge1$ retain the real polynomial
\[
 p_n(y)=\frac{(y+i\pi)^n-(y-i\pi)^n}{2i\pi}.
\]
Define the normalized principal value
\begin{align}
 \mathcal V_n(f)={}&\operatorname{PV}\int_0^1
 p_n(\log x)\frac{f(x)-d_0}{x}\,dx
 +\operatorname{PV}\int_1^\infty
 p_n(\log x)\frac{f(x)}x\,dx.\label{stcut:PV}
\end{align}
The principal value deletes symmetric intervals $(b-\epsilon,b+\epsilon)$
about each pole, with the cutoffs independent. No finite counterterm
is added. There is no pole at $x=1$ under the stated exclusions.

\begin{theorem}[Digamma reconstruction with explicit contour residues]
\label{stcut:reconstruction}
All the principal values in \eqref{stcut:PV} exist, and
\begin{equation}\label{stcut:reconstruct}
 \boxed{d_n^\sigma=\mathcal V_n(f)
 +\mathbf1_{2\mid n}\frac{(-1)^{n/2}\pi^n}{n+1}d_0
 -\sigma i\pi\sum_{b\in\mathcal B}
                  R_b\frac{p_n(\log b)}b.}
\end{equation}
Consequently the exact jump is
\begin{equation}\label{stcut:jump}
 d_n^+-d_n^-=-2\pi i\sum_{b\in\mathcal B}R_b\frac{p_n(\log b)}b.
\end{equation}
The statement permits complex weights; $\mathcal V_n$ is not replaced
by a real-part operation in that case.
\end{theorem}
\begin{proof}
At every positive pole, write
$f(x)=R_b/(x-b)+h_b(x)$ with $h_b$ holomorphic nearby.
The elementary boundary formula is
\[
 \lim_{\eta\downarrow0}\frac1{x-b+\sigma i\eta}
  =\operatorname{PV}\frac1{x-b}-\sigma i\pi\delta_b.
\]
For completeness, subtract the value of a smooth test function at
$b$ in a fixed symmetric neighborhood. The subtracted integral
converges ordinarily, while the constant term is evaluated by an
arctangent and tends to $-\sigma i\pi$; this proves the formula
with the displayed sign and no convention left implicit.

Apply this boundary identity with the smooth local weight
$p_n(\log x)/x$ at each $b$. Away from these finitely many points,
compact convergence is ordinary. Near zero the subtracted numerator
is uniformly $O(x)$, and at infinity the zero-sum digamma asymptotic
is uniformly $O(x^{-1})$. These bounds with the fixed logarithmic
weights justify the two noncompact limits. Taking upper or lower
limits in the complex-shift reconstruction therefore gives
\eqref{stcut:reconstruct}. Subtraction gives \eqref{stcut:jump}.

As an independent sign check, choose integers $N_j\ge0$ such that
$A_j=a_j+N_j>0$ and $a_j+k<0$ for $0\le k<N_j$. The Hurwitz shift
relation gives directly
\[
 d_n^\sigma=\sum_jc_j\gamma_n(A_j)
 -\sum_{b\in\mathcal B}R_b\frac{(\log b+\sigma i\pi)^n}{b}.
\]
Its difference between the two banks is exactly \eqref{stcut:jump}.
This expression is not used to assume the principal-value identity;
it also supplies independent values for verification.
\end{proof}

For example, with $a=-1/2$, $b=1/2$ and weights $1,-1$,
$f(x)=\gamma_0(x-1/2)-\gamma_0(x+1/2)=1/(x-1/2)$.
The formula becomes the elementary, all-index identity
\begin{align}
 &\operatorname{PV}\int_0^1p_n(\log x)
     \frac{(x-1/2)^{-1}+2}{x}\,dx
 +\operatorname{PV}\int_1^\infty
        p_n(\log x)\frac{dx}{x(x-1/2)}\notag\\
 &\qquad=-2\operatorname{Re}(-\log2+i\pi)^n
   +\mathbf1_{2\mid n}\frac{2(-1)^{n/2}\pi^n}{n+1},
 \label{stcut:half-example}
\end{align}
where the use of a real part is confined to this real-weight example.
Its upper or lower Stieltjes value is
$\gamma_n(-1/2\pm i0)-\gamma_n(1/2)
=-2(-\log2\pm i\pi)^n$.

\subsection{The Mellin master and logarithmic cuts}
For a pole at $b>0$, define the higher logarithmic boundary kernel by
\begin{equation}\label{stcut:log-boundary}
 T_{n,b}^{\sigma}
 =\frac1{n+1}\frac{d}{dx}
      \bigl(\Log(x-b+\sigma i0)\bigr)^{n+1}.
\end{equation}
The expression being differentiated is locally integrable. Away from
$b$, the resulting distribution agrees with
$\Log^n(x-b+\sigma i0)/(x-b)$, with
$\Log(x-b+\sigma i0)=\log(b-x)+\sigma i\pi$ for $x<b$.
For $n=0$ this is $\operatorname{PV}(x-b)^{-1}-\sigma i\pi\delta_b$.
For higher $n$, interpreting it as an unchanged real-axis ordinary
integral would discard both logarithmic cuts and local finite parts.

\begin{proposition}[An explicit local cut kernel]
\label{stcut:beta}
For $b>0$, initially $0<\operatorname{Re}v<\operatorname{Re}s<1$,
\begin{align}
 \int_0^\infty x^{v-1}(x-b+\sigma i0)^{-s}\,dx
 &=b^{v-s}\left[e^{-\sigma i\pi s}B(v,1-s)
                        +B(s-v,1-s)\right]\notag\\
 &=e^{\sigma i\pi(v-s)}b^{v-s}
                     \frac{\Gamma(v)\Gamma(s-v)}{\Gamma(s)}.
 \label{stcut:beta-kernel}
\end{align}
The first line exhibits the cut interval $0<x<b$ explicitly. At
$s=1$, continue the combined expression before differentiating:
$(-1)^n\partial_s^n|_{s=1}$ is exactly the Mellin pairing with
$T_{n,b}^{\sigma}$. Its zeroth-index boundary is
\begin{equation}\label{stcut:pole-Mellin}
 \operatorname{PV}\int_0^\infty\frac{x^{v-1}}{x-b}\,dx
 -\sigma i\pi b^{v-1}
 =-\pi b^{v-1}\cot(\pi v)-\sigma i\pi b^{v-1},
 \quad0<\operatorname{Re}v<1.
\end{equation}
\end{proposition}
\begin{proof}
Split at $b$. In the first integral set $x=bt$; in the second set
$x=b/(1-t)$. The two ordinary beta integrals give the first line.
Euler reflection simplifies their sum to the second. Analytic
continuation of $(x-b+\sigma i0)^{-s}$ through $s=1$ is obtained
by one integration by parts against a local smooth test function;
its derivatives give \eqref{stcut:log-boundary}. This establishes
the continued equality and the normalization at every fixed index.
Setting $s=1$ gives the displayed principal-value formula.
\end{proof}

\begin{theorem}[Filtered Mellin transport through the negative shifts]
\label{stcut:Mellin}
Suppose the real nodes and complex weights additionally satisfy
$\sum_jc_ja_j^\ell=0$ for $0\le\ell<r$, where $r\ge1$.
For $0<\operatorname{Re}v<r$, let $\mathcal M_n^\sigma(v)$ be the
Mellin integral of $\sum_jc_j\gamma_n(x+a_j+\sigma i0)$, using
\eqref{stcut:log-boundary} at its positive singularities.
Then
\begin{equation}\label{stcut:Mellin-eq}
 \mathcal M_n^\sigma(v)=(-1)^n\left.\partial_u^n
 \left\{\frac{\Gamma(v)\Gamma(1+u-v)}{\Gamma(1+u)}
       \sum_jc_j\zeta(1+u-v,a_j+\sigma i0)\right\}\right|_{u=0}.
\end{equation}
Every apparent Gamma pole is removed in the combined product before
differentiation. The positive-integer resonant formulas in
\eqref{stmell:integer} therefore remain valid with these boundary
values and this integral prescription.
\end{theorem}
\begin{proof}
Translate every node by the same complex number $z$. All $r$
polynomial cancellations survive, by the binomial formula.
For sufficiently large positive $z$ the positive-shift Mellin theorem
applies. Both sides continue holomorphically while the common shift
avoids the finitely many negative-axis cuts, and hence agree on their
upper and lower continuations to $z=0$.

More explicitly, the Hurwitz shift relation decomposes each integrand
into a positive-node Stieltjes function plus finitely many kernels
\eqref{stcut:log-boundary}. Their local limits exist as just proved.
Near zero the nodes avoid the Hurwitz singularities. At infinity the
original polynomial cancellations give the uniform bound
$O(x^{-r}(1+\log^n x))$, so the tails converge locally uniformly in
$0<\operatorname{Re}v<r$. Thus the limit of the continued integrals
is exactly the specified Mellin pairing, rather than an unchanged
ordinary integral. The local cut kernels also make every residue and
logarithmic contribution explicit. Finally, the Bernoulli cancellations
used at integer Mellin values are polynomial identities in the nodes
and are unaffected by the selected bank, so the existing resonant
coefficient extraction applies without alteration.
\end{proof}

% Scope: this resolves principal-cut lateral transport for the finite
% translation/Mellin families, not arbitrary winding contours, products
% of balanced kernels, or the full general complex-shift research program.
